\section{Entendiendo las matemáticas detrás de la regresión lineal}

Supongamos que tenemos una base de datos hipotética que contiene la información acerca del costo (en unidades de $\$10000$) de varias casas y sus respectivos tamaños (en pies cuadrados $ft^{2}$).


\begin{center}
	\begin{tabular}{|l|l|}\hline
		Tamaño & Costo\\\hline
		1500 & 45\\\hline
		1200 & 38\\\hline
		1700 & 48\\\hline
		800 & 27\\\hline
	\end{tabular}
\end{center}



En este caso, el costo es la variable de salida, mientras que el tamaño es la variable de entrada. 

La entrada y la salida generalmente se denotan por $X$ y $Y$, respectivamente.


En el caso de la regresión lineal, supondremos que el costo $Y$ es una función lineal de tamaño $X$ y para estimar $Y,$ proponemos el modelo \begin{align}
	Y_{e} = \a + \beta X ,
\end{align}
donde $Y_{e}$ es el \emph{valor estimado} de $Y$ con base en nuestra ecuación lineal.


\begin{observacion}
	El propósito de la regresión lineal es encontrar valores $\alpha, \beta$ estadísticamente significativos, que \emph{minimicen} la diferencia entre $Y$ y $Y_{e}.$
\end{observacion}



En el caso de nuestro ejemplo, si encontramos los valores de $\a=2$ y $\beta=0.03,$ entonces la ecuación será
\begin{align}
	Y_{e}= 2 + 0.03 X.
\end{align}


Usando esta ecuación, podemos estimar el costo una casa de cualquier tamaño.  Por ejemplo, para una casa de $900 ft^{2}$, el costo será
\begin{align}
	Y_{e}= 2 + 0.03(900)= 29 .
\end{align}


La siguiente pregunta que nos haremos es como estimar $\a$ y $\beta.$  Para esto usaremos un método llamado suma de  \emph{mínimos cuadrados} para la diferencia entre $Y$ y $Y_{e},$ que representaremos como
\begin{align}
	\ep = Y-Y_{e}.
\end{align}



Nuestro objetivo es minimizar
\begin{align}
	\sum \ep^{2} & = \sum \left( Y-Y_{e} \right)^{2}\\
	&= \sum \left( Y - \left( \a+\beta X \right) \right)^{2}
\end{align}
respecto de los parámetros $\a, \beta.$


Utilizando un poco de cálculo, se puede demostrar que los valores de los parámetros que minimizan la suma anterior son
\begin{align}
	\label{beta}
	\beta &= \dfrac{\sum \left( x-\bar{x} \right)\left( y-\bar{y} \right)}{\sum\left( x-\bar{x} \right)^{2}}\\
	\label{alfa}
	\a &= \bar{y}-\beta\times \bar{x}
\end{align}

\subsection{Demostración de las fórmulas de mínimos cuadrados}

Para encontrar los valores de $\a$ y $\beta$ que minimizan la suma de cuadrados de los errores, definimos la función objetivo:
\begin{align}
	S(\a, \beta) = \sum_{i=1}^{n} \left( Y_i - (\a + \beta X_i) \right)^2
\end{align}

Para minimizar $S$, calculamos las derivadas parciales respecto a $\a$ y $\beta$ e igualamos a cero:

\textbf{Derivada respecto a $\a$:}
\begin{align}
	\frac{\partial S}{\partial \a} &= -2 \sum_{i=1}^{n} \left( Y_i - \a - \beta X_i \right) = 0 \\
	\sum_{i=1}^{n} Y_i - n\a - \beta \sum_{i=1}^{n} X_i &= 0 \\
	n\a &= \sum_{i=1}^{n} Y_i - \beta \sum_{i=1}^{n} X_i \\
	\a &= \bar{Y} - \beta \bar{X}
\end{align}

\textbf{Derivada respecto a $\beta$:}
\begin{align}
	\frac{\partial S}{\partial \beta} &= -2 \sum_{i=1}^{n} X_i \left( Y_i - \a - \beta X_i \right) = 0 \\
	\sum_{i=1}^{n} X_i Y_i - \a \sum_{i=1}^{n} X_i - \beta \sum_{i=1}^{n} X_i^2 &= 0
\end{align}

Sustituyendo $\a = \bar{Y} - \beta \bar{X}$:
\begin{align}
	\sum_{i=1}^{n} X_i Y_i - (\bar{Y} - \beta \bar{X}) \sum_{i=1}^{n} X_i - \beta \sum_{i=1}^{n} X_i^2 &= 0 \\
	\sum_{i=1}^{n} X_i Y_i - \bar{Y} \sum_{i=1}^{n} X_i + \beta \bar{X} \sum_{i=1}^{n} X_i - \beta \sum_{i=1}^{n} X_i^2 &= 0 \\
	\sum_{i=1}^{n} X_i Y_i - n\bar{X}\bar{Y} &= \beta \left( \sum_{i=1}^{n} X_i^2 - n\bar{X}^2 \right) \\
	\beta &= \frac{\sum_{i=1}^{n} X_i Y_i - n\bar{X}\bar{Y}}{\sum_{i=1}^{n} X_i^2 - n\bar{X}^2}
\end{align}

Esta expresión es equivalente a:
\begin{align}
	\beta &= \frac{\sum_{i=1}^{n} (X_i - \bar{X})(Y_i - \bar{Y})}{\sum_{i=1}^{n} (X_i - \bar{X})^2}
\end{align}

\subsection{Descomposición de la suma de cuadrados: SST = SSR + SSD}

Definimos tres cantidades fundamentales:

\begin{itemize}
	\item \textbf{Suma total de cuadrados (SST):}
	\begin{align}
		SST = \sum_{i=1}^{n} (Y_i - \bar{Y})^2
	\end{align}
	Mide la variabilidad total de $Y$.
	
	\item \textbf{Suma de cuadrados de la regresión (SSR):}
	\begin{align}
		SSR = \sum_{i=1}^{n} (\hat{Y}_i - \bar{Y})^2
	\end{align}
	Mide la variabilidad explicada por el modelo.
	
	\item \textbf{Suma de cuadrados de los errores (SSD):}
	\begin{align}
		SSD = \sum_{i=1}^{n} (Y_i - \hat{Y}_i)^2
	\end{align}
	Mide la variabilidad no explicada por el modelo.
\end{itemize}

\begin{teorema}
	$SST = SSR + SSD$
\end{teorema}

\begin{proof}
	Partimos de la identidad:
	\begin{align}
		Y_i - \bar{Y} = (\hat{Y}_i - \bar{Y}) + (Y_i - \hat{Y}_i)
	\end{align}
	
	Elevando al cuadrado y sumando:
	\begin{align}
		\sum_{i=1}^{n} (Y_i - \bar{Y})^2 &= \sum_{i=1}^{n} \left[ (\hat{Y}_i - \bar{Y}) + (Y_i - \hat{Y}_i) \right]^2 \\
		&= \sum_{i=1}^{n} (\hat{Y}_i - \bar{Y})^2 + \sum_{i=1}^{n} (Y_i - \hat{Y}_i)^2 + 2\sum_{i=1}^{n} (\hat{Y}_i - \bar{Y})(Y_i - \hat{Y}_i)
	\end{align}
	
	El término cruzado es cero porque:
	\begin{align}
		\sum_{i=1}^{n} (\hat{Y}_i - \bar{Y})(Y_i - \hat{Y}_i) &= \sum_{i=1}^{n} (\hat{Y}_i - \bar{Y})e_i \\
		&= \sum_{i=1}^{n} \hat{Y}_i e_i - \bar{Y} \sum_{i=1}^{n} e_i
	\end{align}
	
	Por las ecuaciones normales de mínimos cuadrados, sabemos que $\sum_{i=1}^{n} e_i = 0$ y $\sum_{i=1}^{n} X_i e_i = 0$, lo que implica $\sum_{i=1}^{n} \hat{Y}_i e_i = 0$.
	
	Por lo tanto:
	\begin{align}
		SST = SSR + SSD
	\end{align}
\end{proof}

\subsection{Propiedades del coeficiente de determinación $R^2$}

El coeficiente de determinación se define como:
\begin{align}
	R^2 = \frac{SSR}{SST} = 1 - \frac{SSD}{SST}
\end{align}

\begin{observacion}[Propiedades de $R^2$]
	\begin{enumerate}
		\item $0 \leq R^2 \leq 1$
		\item $R^2 = 1$ si y solo si $SSD = 0$, es decir, el modelo explica perfectamente todos los datos.
		\item $R^2 = 0$ si y solo si $SSR = 0$, es decir, el modelo no explica nada de la variabilidad de $Y$.
		\item $R^2$ nunca disminuye al agregar más predictores al modelo (aunque esto no significa que el modelo mejore).
		\item Para regresión lineal simple, $R^2 = r^2$, donde $r$ es el coeficiente de correlación de Pearson entre $X$ y $Y$.
	\end{enumerate}
\end{observacion}

\begin{observacion}[R² ajustado]
	Para compensar el hecho de que $R^2$ siempre aumenta al agregar predictores, se define el $R^2$ ajustado:
	\begin{align}
		R^2_{adj} = 1 - \frac{SSD/(n-p-1)}{SST/(n-1)} = 1 - \frac{n-1}{n-p-1}(1 - R^2)
	\end{align}
	donde $p$ es el número de predictores. El $R^2$ ajustado puede disminuir si se agrega un predictor que no mejora suficientemente el modelo.
\end{observacion}

% \part{NO BORRAR}
% \section{Recursos}
% \subsection{analisisPredictivo}
% El repositorio con los scripts de \texttt{Python 3} de esta presentación los puede encontrar en \href{https://github.com/julihocc/ulsaPye/tree/master/analisisPredictivo}{https://github.com/julihocc/ulsaPye/tree/master/analisisPredictivo}
% 
% [t, ]
% \frametitle{Referencias}
% \nocite{*}
% \betaibliographystyle{amsalpha}
% \betaibliography{ulsaPye}
% 